\section{Corrections to the literature}
\label{sec:correc}

The three statements in this section are about published work rather than about
our constructions, and we have tried to state them so that a reader can check
each against the primary source without our code.

\subsection{The announced \texorpdfstring{$(42,5)$}{(42,5)} does not follow from the method cited for it}

On p.~450 of \cite{jsww1976} one reads:

\begin{quote}
``All that is necessary to reduce the degree to 5 is the Skolem substitution
method (cf.\ [3], p.~263). However, this procedure increases the number of
variables (to 42 when applied to (1)).''
\end{quote}

Their reference [3] is \cite{davis1973}, and on p.~263 the method is given in
full. It permits exactly the substitutions
\[
  z_j=y_iy_k,\qquad z_j=y_i^2,\qquad z_j=xy_i,\qquad z_j=x^2 ,
\]
applied successively, so that the admissible names are precisely the
\emph{degree-2 monomials} in the variables, including those already introduced.
Composite subexpressions are not admissible: of the twenty names in our
construction, four---among them $(k+1)^3(k+2)$, $cu+x$ and $a+u^2(u^2-a)$---are
not monomials and would not be available.

This makes the minimum computable exactly rather than heuristically. The
candidate names are the divisors of the monomials occurring in the system (a
name dividing none of them is useless), and one enumerates all partitions of
exponents subject to a name of degree greater than $2$ being built from two
already-available parts, which is what ``successive substitutions'' means. On
the system~(1) of \cite{jsww1976} the exact minimum is $25$ names, so the method
applied to that system yields
\[
  26+25 \;=\; 51 \text{ variables at degree } 5 ,
\]
that is, $(51,5)$, not $(42,5)$. Reaching $42$ by that route would require $16$
names, and no selection of $16$ monomials flattens the system.

We are careful about what this does and does not show. It does not show that no
$(42,5)$ generator exists; it shows that it does not follow from the procedure
cited for it. Our $(44,5)$, which uses composite names, beats that procedure by
seven variables. This is a statement about an attribution, and it is the only
one of our results that we would describe as surprising.

\subsection{Two conflations of units}

\paragraph{The figure $1.638\cdot10^{45}$ is not the degree of a prime-representing polynomial.}
It is the degree of the \emph{universal} pair $(9,\cdot)_\N$ of \cite{jones1982}:
a bound valid for every Diophantine set, not a polynomial for the primes. The
prime-representing polynomial in $10$ variables that Matiyasevich actually
constructed \cite{matiyasevich1981} has degree greater than $6000$, a figure
confirmed by its formalizers \cite{pak2022ten}. The two are routinely quoted as
one, with the $10^{45}$ introduced by a ``hence'' from the nine-unknowns
theorem; that is a different polynomial, and not a constructed one.

\paragraph{The pair $(58,4)$ is neither about the primes nor about a generator.}
It is a universal pair, and a pair for a \emph{representation}. Instantiated for
the primes it would give $59$ variables, and as a generator its degree would be
$1+2\cdot4=9$; it is dominated on both axes. We note in passing that no
universal pair can have degree $\le2$, since quadratic Diophantine equations are
decidable.

\subsection{A consistency check that catches such conflations}

By Proposition~\ref{prop:gen}, the degree of a generator is odd. A quoted figure
that is even is a representation degree, not a generator degree; a figure that
is odd and equals $1+2d$ for a plausible $d$ is a generator degree. Applied to
$13697=1+2\cdot6848$ this confirms the reading of the $(12,13697)$
of \cite{jsww1976}, and applied to $(58,4)$ it immediately flags the unit
mismatch. We recommend it as a routine check when citing these figures.
